

% \subsection{Recursive Reasoning}
% \label{sec:recursion}

% The recursive reasoning module serves as the central controller of \mot, 
% coordinating temporal decomposition, toolkit-driven retrieval, sufficiency verification, and final answer synthesis. 
% It acts as the reasoning router that dynamically determines whether to recall prior experiences, 
% select appropriate temporal toolkits, perform hybrid retrieval, or refine sub-questions. 
% By iteratively executing and verifying each sub-question, the framework ensures both semantic fidelity and temporal consistency throughout the reasoning process.



% \myparagraph{Question tree construction}
% Given a temporal question $Q$ and its classified temporal type $\textsf{Type}(Q)$, 
% the reasoning process begins by constructing a hierarchical decomposition tree $\mathcal{T}_Q$.  
% Each node in this tree represents a sub-question capturing a specific temporal constraint or relation.  
% To enhance decomposition robustness, \mot\ retrieves decomposition exemplars from the experience pool $\mathcal{E}_{pool}$ and integrates them into the LLM prompt:
% \[
% \mathcal{E}_{\text{decomp}} = \textsf{GetDecompExperience}(\mathcal{E}_{pool}, Q, W_{\text{exp}}, \textsf{Type}(Q)),
% \]
% \[
% \mathcal{T}_Q = \textsf{Prompt}_{\text{LLMDecompose}}(Q, \mathcal{E}_{\text{decomp}}, \textsf{Type}(Q)).
% \]
% This experience-guided decomposition enables the model to adaptively construct more faithful and temporally aligned query structures.  
% Each node $q_i \in \mathcal{T}_Q$ is represented as:
% \[
% q_i = (\textit{idx},\, \textit{question\_text},\, \textit{operator},\, \textit{constraint},\, \textit{children}),
% \]
% where \textit{operator} specifies the temporal relation (\eg \textsc{before}, \textsc{after}, \textsc{within-range}), 
% and \textit{constraint} denotes the time condition.  
% For any parent–child pair $(q_i, q_j)$, the temporal order satisfies $t(q_i) \le t(q_j)$, 
% ensuring the reasoning process respects temporal monotonicity.

% \myparagraphunderline{Indicator extraction}
% From the constructed tree, the system extracts a set of sub-question indicators $\{\Ind_i\}$ 
% that encode the semantic and temporal constraints of each node:
% \[
% \{\Ind_i\} = \textsf{ExtractIndicators}(\mathcal{T}_Q),
% \]
% where each indicator $\Ind_i = \langle x?, R, y?, \tau; \mathcal{C}_{time}\rangle$ specifies the entities, relation, and temporal scope for reasoning.  
% The predicted reasoning depth $D_{\text{pred}}(q_i)$ is estimated as the maximum distance between each sub-question’s predicted answer and its topic entities, 
% serving as a heuristic for recursion control.








% % \subsection{Hierarchical Reasoning}
% % \label{sec:recursion}

% % \myparagraph{Hierarchical execution control}
% % Given the reasoning tree $\mathcal{T}_Q = \{q_1, q_2, \dots, q_n\}$ and their indicators $\{\Ind_i\}$, 
% % the hierarchical reasoning loop traverses the tree in a top–down manner, 
% % propagating verified temporal constraints to child nodes and aggregating resolved answers upward.  
% % % Each sub-question $q_i$ is processed according to a unified control policy:
% % At each step, \mot first to examin if the exsiting retrieved knowledge is able to answer the given question. if not it will invoke temporal toolkits for new templora evidence exploration:
% % \[
% % \textsf{Execute}(q_i) =
% % \begin{cases}
% % \textsf{MemoryLookup}(q_i, \mathcal{E}_{pool}), & \text{if matches},\\
% % \textsf{ToolkitSelection}(q_i, \textsf{Type}(Q), \mathcal{T}_{\theta}), & \text{otherwise.}
% % \end{cases}
% % \]


% \myparagraph{Hierarchical execution control}
% Given the reasoning tree $\mathcal{T}_Q = \{q_1, q_2, \dots, q_n\}$ and their indicators $\{\Ind_i\}$, 
% the hierarchical reasoning loop traverses the tree in a top–down manner, 
% propagating verified temporal constraints to child nodes and aggregating resolved answers upward.  
% At each step, \mot\ first examines whether existing retrieved knowledge can sufficiently answer the sub-question; 
% if not, it invokes temporal toolkits for new evidence exploration:
% \[
% \textsf{Execute}(q_i) =
% \begin{cases}
% \textsf{MemoryLookup}(q_i, \mathcal{E}_{pool}), & \text{if match},\\[4pt]
% \textsf{ToolkitSelection}(q_i, \textsf{Type}(Q), \mathcal{T}_{\theta}), & \text{otherwise.}
% \end{cases}
% \]


% % \myparagraphunderline{Memory Lookup}
% % % \textsf{MemoryLookup} searches the experience pool $\mathcal{E}_{pool}$ for similar sub-questions 
% % % and directly reuses their verified answers if temporal constraints align.  
% % % If no matching experience or sufficiency failure occurs, 
% % % \textsf{ToolkitSelection} is triggered to identify the most suitable retrieval tool for $q_i$ 
% % % based on its indicator, type and toolkit library $\mathcal{T}_{\theta}$.
% % The system first queries the experience pool $\mathcal{E}_{pool}$ for matching trajectories using the indicator $\Ind_i$.  
% % If a stored path $P_{\text{exp}}$ satisfies $\Ind_i$ and pass the corresponding answer and reasoning trace are directly reused for the following suffuicient evluation, .  
% % Otherwise or sufficent evaluation fails, \mot proceeds to select temporal toolkits.




% \myparagraphunderline{Memory lookup.}
% The system first queries the experience pool $\mathcal{E}_{pool}$ for similar sub-question trajectories.  
% If a stored path $P_{\text{exp}}$ satisfies $\Ind_i$, the corresponding reasoning trace and answer are directly reused for sufficiency evaluation.  
% Otherwise, or sufficiency test fails, \mot\ proceeds to the temporal toolkit selection stage.


% % \myparagraphunderline{Temporal toolkit selection}
% % When the memory lookup procedure fails,\mot utilzied the temporal toolkit selection for the detection the new evidence in the section 4.3.  Toolkit selection enables \mot\ to adaptively invoke specialized retrieval strategies for each sub-question.  
% % Instead of relying on a single fixed retriever, \mot\ employs a set of temporal toolkits, 
% % each designed for distinct temporal reasoning purposes such as event ordering, range detection, and interval comparison.  




% \myparagraphunderline{Temporal toolkit selection}
% When memory lookup fails, \mot\ activates temporal toolkit selection to explore new evidence sources, detailed in Section~\ref{sec:retrieval}.  
% Toolkit selection enables \mot\ to adaptively invoke specialized temporal reasoning tools instead of relying on a fixed retriever.  
% Each toolkit is designed for a specific reasoning purpose, such as event ordering, range detection, or interval comparison.
% Following the exemplar-based prompting strategy described in Section~\ref{sec:method:grounding}, 
% \mot\ first retrieves relevant toolkit exemplars from the experience pool:
% \[
% \mathcal{E}_{\text{tool}} = 
% \textsf{GetToolkitExperience}(\mathcal{E}_{pool}, q_i, \Ind_i, \textsf{Type}(Q), W_{\text{exp}}),
% \]
% and constructs a prompt for toolkit selection:
% \[
% \mathcal{T}_i = 
% \textsf{Prompt}_{\text{ToolkitSelectBasedExp}}
% (\mathcal{T}_\theta,\, \Ind_i,\, q_i,\, \mathcal{E}_{\text{tool}}).
% \]
% The LLM may select multiple candidate tool requests 
% $\mathcal{R}_i = \{r^{(1)}, r^{(2)}, \dots, r^{(m)}\}$,  
% each corresponding to a different reasoning strategy.  
% All requests are executed in parallel, producing intermediate results 
% $\mathcal{C} = \{c^{(1)}, c^{(2)}, \dots, c^{(m)}\}$, 
% which are evaluated in the sufficiency testing step.






% % similart to the previouse expeienc enheanced prompt in section 3.1, \mot first extract the closed examplers form the expeience pool $\mathcal{E}_{pool}$. Previouse sufficiful expeience will be first detection to enhance LLM selection ability and decresed the hallucination caused by variouse avalilable toolkit selection: 
% % \[
% % \mathcal{E}_{\text{tool}} = \textsf{GetToolkitExperience}(\mathcal{E}_{pool}, q_i, \Ind_i,\text{Type}(Q), W_{\text{exp}}),
% % \]
% % and constructs a prompt for toolkit selection using these exemplars:
% % \[
% % \mathcal{T}_i = 
% % \textsf{Prompt}_{\text{ToolkitSelectBasedExp}}
% % (\mathcal{T}_\theta,\, \Ind_i,\, q_i,\, \mathcal{E}_{\text{tool}}).
% % \]
% % The LLM may select multiple candidate tool requests 
% % $\mathcal{R}_i = \{r^{(1)}, r^{(2)}, \dots, r^{(m)}\}$,  
% % each representing a different reasoning strategy (\eg event ordering, range comparison, or timeline extraction).  
% % All requests are executed in parallel throught the hybired retrieval propsoed in Section 4.3, producing intermediate results 
% % $\mathcal{C} = \{c^{(1)}, c^{(2)}, \dots, c^{(m)}\}$. All the candidate results will be examed in the following suffcient check procedure.

% % \myparagraph{Question answering}
% % after each subquestion obatined candaide tempral path, the \mot will exam to evlauated if it is suuficent answer the current answer, and up to the overall question, if no, a tree update will be processed. 

% % \myparagraphunderline{Subquestion evluation}
% % after each subquestion obatined candaide tempral path, the \mot will prompt LLm to evaluate if the current result is sufficient answer the question and then generate a answer for the given question in a signle LLM call. if multiple diffrent answer obatined and pass the suffcient checking throughe the different toolkit retreval, a debated-vote method is utilziaed to select the best result.
% % The node $q_i$ is then labeled as solved.





% % \myparagraphunderline{Global sufficiency and answer generation}
% % After all sub-questions are processed, \mot performs a global sufficiency test over the entire reasoning tree to ensure temporal and logical coherence. after the sufficent pass, it move to the next answer generation pahse.
% % in oreder to reduce the hallucanation caused by the massive irrevlent knowledge and cannot distingush the temporal informatiom.
% % All the previouse temporal path will be summaized into a consine tmpoeral reasoning path, that align with the indicator of orginal question. The summaried temoporal reasoning will be used for prompt LLM to generate an answer. 
% % The exprience of while trajtory and toolkit selection experience will be update to the $\mathcal{E}_{pool}$.

% % \myparagraphunderline{Tree update and termination.}
% % When sufficiency verification fails, \mot adaptively decides among keep retreval, refinement query and further decomposition based on the current knwoledge solving sitation and recursion depth. 
% % \[
% % \{q_{i*}\} = 
% % \textsf{RefineOrDecompose}(q_i, \Ind_i, \mathcal{G}_Q, D_{\text{cur}}, D_{\max}),
% % \]
% % The new child node will be updated to $\mathcal{T}_Q$.  
% % The recursion terminates when all leaf nodes are marked as solved or when the maximum depth $D_{\max}$ is reached, 
% % after which the verified paths are concatenated to form the final reasoning chain $\mathcal{P}^{\text{final}}$ for answer generation.


% \myparagraph{Question answering}
% After each sub-question $q_i$ obtains candidate temporal results, 
% \mot\ evaluates whether these results sufficiently answer the local query and contribute to the overall question. 
% This stage combines evidence evaluation, global reasoning synthesis, and hierarchical update for unresolved nodes.

% \myparagraphunderline{Sub-question evaluation}
% For each node in the reasoning tree, the LLM assesses retrieved temporal paths for both semantic relevance and temporal validity. 
% A sub-question is considered sufficient when one or more candidate paths provide consistent evidence aligned with its indicator. 
% If multiple valid results exist, a lightweight debate–vote strategy aggregates them into a single coherent answer. 
% The solved node and its reasoning trace are recorded in the experience pool $\mathcal{E}_{pool}$ for future reuse.

% \myparagraphunderline{Tree update and termination}
% If a sub-question fails sufficiency verification, \mot\ adaptively decides whether to retry retrieval, refine the query, or perform further decomposition based on its retrieved information. 
% Newly generated nodes are appended to the tree and processed hierarchically until all branches are solved or the maximum depth $D_{\max}$ is reached. 
% Finally, all verified reasoning paths are consolidated into the final evidence for answer generation.

% \myparagraphunderline{Global sufficiency and answer synthesis}
% After all sub-questions are processed, \mot\ checks the temporal and logical coherence of the entire reasoning tree. 
% When global sufficiency is achieved, verified sub-paths are merged into a unified temporal reasoning path that summarizes the full evidence flow. 
% This concise reasoning trace is then used to prompt the final LLM for answer generation, ensuring interpretability and temporal consistency.



% % \myparagraph{Question answering.}
% % Once each sub-question $q_i$ obtains its candidate temporal results, 
% % \mot\ proceeds to evaluate whether these results are sufficient to answer both the local sub-question and, hierarchically, the global question. 
% % This stage integrates evidence evaluation, global reasoning consistency, and dynamic tree update.

% % \myparagraphunderline{Sub-question evaluation.}
% % For each node in the reasoning tree, \mot\ first examines the temporal paths retrieved from the corresponding toolkit calls. 
% % These paths are assessed by the LLM for both semantic relevance and temporal validity. 
% % A sub-question is considered sufficient when at least one candidate path provides clear and temporally consistent evidence that supports its indicator. 
% % If multiple toolkits yield valid results, \mot\ aggregates them through a lightweight debate–vote mechanism, 
% % where each candidate explanation is compared and the most coherent one is selected as the final local answer.  
% % The node is then marked as \textit{solved}, and both its reasoning trajectory and selected toolkit are recorded in the experience pool $\mathcal{E}_{pool}$ for future reuse.  
% % This experience-grounded reuse mechanism allows \mot\ to gradually improve its efficiency and reliability with each new question.

% % \myparagraphunderline{Global sufficiency and answer synthesis.}
% % After all sub-questions have been executed, \mot\ performs a global consistency verification over the entire reasoning tree. 
% % This step ensures that all local sub-answers are temporally ordered, logically coherent, and aligned with the global temporal constraint of the original question. 
% % When the reasoning tree achieves global sufficiency, 
% % the verified sub-paths are integrated into a single temporal reasoning chain that captures the complete chronological progression of evidence.  
% % To mitigate potential hallucination or redundancy, \mot\ summarizes this reasoning chain into a concise textual explanation that highlights key temporal transitions and their supporting facts.  
% % This summarized reasoning trace is then fed into the final LLM call to generate the natural-language answer for the original question.  
% % During this process, \mot\ not only produces an interpretable output but also ensures that every reasoning step is traceable and temporally grounded.  
% % The entire trajectory—including the selected tools, verified evidence, and generated answer—is subsequently stored in the memory layer for continual learning.

% % \myparagraphunderline{Tree update and termination.}
% % When a sub-question fails sufficiency verification, \mot\ adaptively determines the next step according to its current reasoning depth and available contextual information. 
% % If sufficient evidence may still exist in the temporal graph, the system re-invokes the retrieval pipeline for refinement; 
% % otherwise, the sub-question is further decomposed into smaller queries to explore finer-grained temporal relations.  
% % Newly generated sub-nodes are appended to the reasoning tree, and the process continues hierarchically until either all nodes are solved or the maximum reasoning depth $D_{\max}$ is reached.  
% % Once all leaf nodes are resolved or pruned, the verified temporal reasoning paths are concatenated to form the final reasoning chain, 
% % which is then passed to the answer generation stage.  
% % Through this hierarchical loop of evaluation, refinement, and synthesis, 
% % \mot\ achieves an adaptive reasoning process that unifies temporal understanding, factual grounding, and dynamic memory utilization.


\vspace{-4mm}
\subsection{Tree of Time — Hierarchical Temporal Reasoning}
\label{sec:recursion}

% The hierarchical reasoning module serves as the central controller of \mot, 
% coordinating temporal decomposition, toolkit-driven execution, sufficiency verification, and final answer synthesis. 
% It acts as the reasoning router that dynamically determines whether to recall prior experiences, 
% select temporal toolkits, refine unresolved nodes, or finalize the reasoning chain. 
% By iteratively executing and validating sub-questions, 
% the framework maintains both semantic fidelity and temporal consistency throughout the reasoning process.

% After temporal grounding, \mot\ enters the \emph{hierarchical reasoning} stage, 
% where grounded temporal structures are progressively transformed into executable reasoning steps.  
% This module serves as the central controller of the framework, 
% coordinating hierarchical decomposition, toolkit execution, sufficiency verification, and answer synthesis. 
% It dynamically decides when to recall prior experience, invoke toolkits, or refine unresolved sub-questions, 
% ensuring semantic fidelity and temporal consistency throughout the reasoning process.

% After temporal grounding, \mot\ proceeds to the core reasoning stage, 
% where grounded temporal representations are transformed into executable reasoning procedures.  
% This stage coordinates hierarchical decomposition, toolkit-driven execution, sufficiency verification, and final answer synthesis.  
% It acts as the central controller that dynamically determines operators.
% when to recall prior experience, 
% invoke temporal toolkits, refine unresolved sub-questions, or finalize the reasoning chain.  
% Through iterative execution and validation, \mot\ maintains both semantic fidelity and temporal coherence across the reasoning process.



% After temporal grounding, \mot\ enters the {hierarchical reasoning} stage, 
% where grounded temporal structures are progressively transformed into executable reasoning steps.  
% This stage serves as the central controller of the framework, 
% coordinating hierarchical decomposition, toolkit execution, sufficiency verification, and answer synthesis. 
% It adaptively balances memory recall, toolkit invocation, and sub-question refinement throughout the reasoning.
After temporal grounding, \mot\ enters the Tree of Time (ToT) stage,
which transforms grounded temporal structures into a hierarchical reasoning tree.
Serving as the framework’s central controller, it coordinates decomposition, toolkit execution, sufficiency testing, and answer synthesis.
Guided by the global plan, the model adaptively alternates between recalling experience, invoking operator-specific toolkits, and refining sub-questions to maintain temporal and semantic consistency.
% Due to the space limitation, the pseudo-code of this section is shown in Appendix \ref{algorithm:analysis}.
Due to the space limitation, 
the pseudo-code of the ToT is detailed in Algorithm \ref{algorithm:controller} of Appendix \ref{appendix:alg}.
% and the prompt template is presented in Appendix \ref{appendix:prompt}.

\myparagraph{Question tree construction}
Given a temporal question $Q$ and its classified temporal type $\textsf{Type}(Q)$, 
\mot\ begins by constructing a hierarchical decomposition tree $\mathcal{T}_Q$.  
Each node represents a sub-question with a specific temporal relation or constraint.  
To enhance robustness, decomposition exemplars are retrieved from the experience pool $\mathcal{E}_{pool}$ and integrated into the LLM prompt:
\[
\mathcal{E}_{\text{decomp}} = \textsf{GetDecompExp}(\mathcal{E}_{pool}, Q, W_{\text{exp}}, \textsf{Type}(Q)),
\]
\[
\mathcal{T}_Q = \textsf{Prompt}_{\text{Decompose}}(Q, \mathcal{E}_{\text{decomp}}, \textsf{Type}(Q)).
\]
This experience-guided process enables the system to build a type-aligned reasoning tree that reflects temporal dependencies among sub-questions. 
For any parent–child pair $(q_i, q_j)$, the temporal order satisfies $t(q_i) \le t(q_j)$, ensuring temporal monotonicity.

\myparagraphunderline{Indicator extraction}
From $\mathcal{T}_Q$, the system extracts a set of sub-question indicators $\{\Ind_i\}$, i.e., $\Ind_i = \langle x?, R, y?, \mathcal{C}_{time}\rangle$,
each encoding the entities, relations, and temporal constraints required for reasoning.  
Based on this, a predicted depth $D_{\text{pred}}(q_i)$ is calculated, defined as the maximum distance between the predicted answer and each topic entity.
For instance, in Figure~\ref{sec:intro}(d), the predicted depth of the overall indicator is 2.

% The predicted reasoning depth $D_{\text{pred}}(q_i)$ is estimated by it then for the complexity of each node, and is later used for recursion control. For the example, in Figure \ref{sec:intro}(d), the predict depth of the overall indicator is 2.

\myparagraph{Hierarchical execution control}
The reasoning tree is traversed in a top-down manner.  
At each node, \mot\ first checks whether a similar reasoning trace exists in the experience pool for direct reuse. 
If no match or evidence insufficient, the system proceeds to dynamic toolkit selection to discover new evidence.

\myparagraphunderline{Experience-guided toolkit selection}
When memory lookup fails, \mot\ retrieves prior examples of successful toolkit usage to enhance selection accuracy:
\[
\mathcal{E}_{\text{tool}} = \textsf{GetToolkitExp}(\mathcal{E}_{pool}, q_i, \Ind_i, \textsf{Type}(Q), W_{\text{exp}}),
\]
and constructs a context-enriched prompt:
$
\mathcal{T}_i = 
\textsf{Prompt}_{\text{ToolkitSelect}}$ $
(\mathcal{T}_\theta,\, \Ind_i,\, q_i,\, \mathcal{E}_{\text{tool}}).
$
The LLM may recommend multiple toolkits simultaneously, 
each representing a distinct retrieval or reasoning strategy (\eg event ordering, interval comparison, or timeline construction).  
All selected tools are executed in parallel, and their candidate results are passed to the sufficiency evaluation stage.

\myparagraph{Question answering}
After each sub-question $q_i$ obtains its candidate temporal results, 
\mot\ evaluates whether these results sufficiently answer the local query and contribute to the overall question. 
This stage integrates evidence evaluation, global reasoning synthesis, and hierarchical update for unresolved nodes.

\myparagraphunderline{Sub-question evaluation}
For each node in the reasoning tree, the LLM assesses retrieved temporal paths for both semantic relevance and temporal validity. 
A sub-question is considered sufficient when one or more candidate paths provide consistent evidence aligned with its indicator. 
If multiple valid results exist, a lightweight debate–vote strategy aggregates them into a single coherent answer. 
The solved node and its reasoning trace are recorded in the experience pool $\mathcal{E}_{pool}$ for future reuse.



\myparagraphunderline{Global sufficiency and answer synthesis}
After all sub-questions are processed, \mot\ checks the temporal and logical coherence of the entire reasoning tree. 
When global sufficiency is achieved, verified sub-paths are merged into a unified temporal reasoning chain that summarizes the full evidence flow. 
This concise reasoning trace is then used to prompt the final LLM for answer generation, ensuring interpretability and temporal consistency.

\myparagraphunderline{Tree update and termination}
If a sub-question fails sufficiency verification, \mot\ adaptively decides whether to retry retrieval, refine the query, or perform further decomposition based on its current reasoning. 
Newly generated nodes are appended to the tree and processed recursively until all branches are solved or the maximum depth $D_{\max}$ is reached. 
Finally, verified reasoning paths are consolidated into the final evidence chain for answer generation.
\vspace{-2mm}
\subsection{Temporal Evidence Retrieval and Pruning}
\label{sec:retrieval}
As discussed in Section~\ref{sec:intro}, identifying reasoning paths that connect all topic entities is crucial for deriving accurate answers.  These paths function as interpretable chains of thought.
% These paths function as interpretable chains of thought, revealing not only the final answer but also the intermediate reasoning steps, a property we refer to as \textbf{interpretability}.  
In this section, we operationalize the selected toolkits to perform hybrid retrieval and pruning, 
balancing {efficiency} and {accuracy} while ensuring both high temporal alignment and broad semantic coverage.
% Due to the space limitation, 
The pseudo-code of this section is detailed in Algorithm \ref{algorithm:execution} (Appendix\ref{appendix:alg}).

